% Page 045 — page imprimée 41
% Contenu seul : le préambule, l'en-tête et la pagination sont gérés par meyrueis.tex

qui se mit à notre service et nous épaula avec beaucoup de dévouement. Les jours de
veillées ou de réunions, Irène pouvait m’accompagner. Au retour, dans la nuit, je
ramenai Jeanne chez elle.

Il y eut, au cours de notre séjour, ici et là, quelques soins médicaux à donner. Bien
plus tard, venant d’Afrique du Nord en vacances et de passage dans la région, tel ou
tel ancien paroissien de Meyrueis ou de Gatuzières nous disait : « Vous souvenez-
vous, Madame Gounelle vous avez soigné ma fille », « vous avez guéri la grand
mère », » « vous avez avant notre départ donné le landau d’André ce qui nous a
soulagés » « Quand ma mère, alitée depuis longtemps, souffrait de nombreuses
escarres, vous avez arrangé son lit et amélioré sa position en sorte que ses
souffrances furent considérablement allégées mais une heure plus tard elle
recommença à souffrir et me disais : Si notre dame était là ! »
Ce côté de notre mission restait bien réel. Mais au Pompidou, évidemment, sans
docteur et sans bonnes sœurs la trace s’était montrée bien plus profonde.
C’est ainsi que, sans éclat particulier, dans la fidélité aux petites choses, se déroula un
ministère paroissial normal, sans heurts et sans secousses, pendant trois ans.

\subsection*{Notre concierge Monsieur Avesque (dit Lou Rascle)}

Un autre trait particulier du Meyrueis de 1930 c’est la personnalité de notre
concierge. Il s’appelait monsieur Avesque. Grand et original, fin parfois dans ses
réflexions, obstiné dans son comportement, vulgaire trop souvent dans son langage, il
avait le don d’imposer sa manière de faire, sans nuance. Très âgé, il avait dû, dans sa
jeunesse être un vrai colosse. Courbé par les ans, il restait encore très vigoureux.
Chaque semaine, il passait une journée entière, pour une rémunération dérisoire, à
nettoyer ce temple immense, briquer ses pavés disjoints, enlever la poussière de ses
bancs tous numérotés, garnir l’hiver les trois grands poêles, bien incapables d’ailleurs
de fournir une chaleur suffisante. Aussi chacun portait son chauffe-pied pour le
service du dimanche, qu’il avait préparé aussi bien que possible.

Après avoir allumé les feux ( il disait éclairé le feu ), le moment venu, il s’attaquait à
la lourde cloche, au son grave et profond, qui se répercutait au loin dans les vallées,
trois fois dans la matinée. Le dernier coup indiquait le début de l’office. Très souvent,
l’église catholique annonçait sa messe à la même heure. M. Avesque n’admettait pas
que nos « frères séparés » (ainsi les désignait-il) mettent leur cloche en branle tant
qu’il n’avait pas terminé. S’il avait commencé le premier, on devait respecter sa
priorité. Sinon il n’arrêtait pas de tirer sur sa corde. Cela aurait pu durer une demi-
heure, ou davantage, il ne cédait pas souvent, impuissant devant cette guerre des
cloches, j’ai été contraint de retarder « ma cérémonie » (sic) car il fallait bien, que le
silence se fit.

\subsection*{Vie de l’église}

L’école du dimanche précédant le culte réunissait de nombreux élèves, pas toujours
très disciplinés. Je tenais beaucoup à ce qu’il y règne une certaine liberté et que cette
école ne ressemble pas, cet égard, à l’école de chaque jour. Nous avions plusieurs
groupes avant l’assemblée générale. Après le rassemblement, quand M. Avesque
s’apercevait qu’un garçon chahutait un peu trop à son gré, ne me trouvant pas assez
sévère, il faisait chauffer son tisonnier à blanc, puis, en pleine leçon il l’approchait,
